\begin{center}
{\Large\bfseries Frobenius traces, marked points and residual chains}\\[4pt]
Edition 0.64 \quad 24 November 2026
\end{center}

For each rank-three host in Case V, the first chapter constructs a single
Frobenius element at which every eligible point has trace different from
the base point. A uniform bound controls the sum of the valuations of these
differences. The second chapter proves irreducibility on the open
subgroups and on the full cyclotomic kernel, and identifies the arithmetic
input still needed to obtain automorphy at the same mark.

In Case III the third chapter constructs a tame tower with the prescribed
label and four distinct residual characters. The fourth chapter connects
the restricted lattice to an integral seed by four polynomial families on
every sufficiently deep later tower, under the stated seed, fitness and
cohomology hypotheses.

\begingroup\emergencystretch=2em
\section{A single separating Frobenius and a uniform trace bound for Case V}
\label{sec:case-v-host-number}

% BEGIN VUN1
\paragraph{VUN1. Restricting a host point to the Gaussian field.}
\label{claim:VUN1}
Let $K$ be any CM field satisfying $\mathrm{IsQi}$: it has degree
two over $\mathbb Q$ and contains a root of $X^2+1$. Work with the
source data $\Omega,\ell,c,N,e,\psi,o_1,S,(A_v),F,\iota,g_1$
in $\mathrm{HostNumber}_s$, including $\mathrm{HostData}$,
$\mathrm{IsHostFrame}$ and $F.\mathrm{Marked}(g_1)$. Write
$O=c.O$, $k=O/\mathfrak m_O$, $\Delta=G_{\mathbb Q}$,
$\pi:\Delta\twoheadrightarrow\Gamma=G_S$ and $H=\ker\eta$ in
$\Gamma$. Put $L=\overline{\mathbb Q}_\ell$ and write $\varepsilon$
for the cyclotomic character. The frame identities are
$\eta\pi=\eta_4$ and $\Psi\pi=\varepsilon^e\psi_{\rm Gal}$.

Take any marked characteristic-zero point $r$ of type I satisfying
the source $B'_s(K,\iota)$ conditions. Its automorphy clause gives
an embedding $j:K\hookrightarrow\Omega$, a representation
$\rho:G_K\to\mathrm{GL}_3(L)$ with matrices equal to those of
$r\pi$ on $G_K$, and the signed automorphy assertion for
$\rho\varepsilon^{-1}$ (HSR8). Both possible $j$ have image
$\mathbb Q(i)$ inside $\Omega$. Thus restriction identifies
$G_K$ with $\ker\eta_4\subset\Delta$, and its composite with
$\pi$ maps onto $H$. Indeed every lift to $\Delta$ of $h\in H$
has trivial $\eta_4$. A $\rho$-invariant subspace is invariant
under all the matrices on $H$, so type I gives irreducibility
over $L$. For fixed $j$, any isomorphism between two such
restrictions also intertwines all matrices on $H$.

The first conjunct of the supplied $\mathrm{IsAutCsdA}$ is
$\mathrm{AutCsd0}$. By HSR8 it consists exactly of a regular
algebraic cusp, its conjugate self-duality and the same attached
system as in $\mathrm{IsAutCsd}$, with twist $\varepsilon^{-1}$.
The matrices of $r$ are continuous, since a torus conjugate comes
from an adically continuous marked point over a finite coefficient
DVR (HSR1). Thus the matrices of $\rho$ and their inverses are
continuous, the latter also by $\rho(g)^{-1}=\rho(g^{-1})$.
For the $\mathrm{Adm3}$ branch of RNC16 we have proved the
irreducibility and automorphy conditions and the continuity field
of the level. We next prove its two local conditions at the
unchanged level $N$.
% END VUN1

% BEGIN VUN2
\paragraph{VUN2. The rank-three level on inertia away from $\ell$.}
\label{claim:VUN2}
Put $d_0=1,d_1=0,d_2=-1$, and on $G_K$ put
$\xi_i=(\psi_{\rm Gal}|_{G_K})^{d_i}$, embedded in $L^\times$.
These three global characters kill the group fixing $\mu_N$.
For every finite place $w\nmid\ell$ of $K$ there is an extension
of $w$ to $\Omega$ whose inertia under $\rho$ is simultaneously
diagonal in one basis, with these three characters.

Here are the local comparisons. Work first at the chosen $A_v$ over
a rational prime $p\ne\ell$ belonging to $S$, in an integral
model $r_D$ of the characteristic-zero point, with coefficient
DVR $D$. Its image over $L$ is a fixed torus conjugate of $r$.
On sign-positive inertia $I^+=I_v\cap\ker(\eta\pi)$ its
multiplier is one, and the base matrix is
$T(\Psi)=\operatorname{diag}(\Psi,1,\Psi^{-1})$.
At $v\in S_b$ the actual diagonal local basis of the integral
point gives exactly $T(\Psi)$ on $I^+$. At
$v\notin T_b=\{v_\ell\}\cup S_b$ both $r_D$ and $T(\Psi)$
kill $N_v=I^+\cap\ker\bar\Psi$ (CVI2). Their restrictions
therefore descend to matrices $U(t),V(t)$ on
$C=I^+/N_v\subseteq k^\times$, where $V=T(\Psi)$.
They agree upon reduction modulo $\mathfrak m_D$, and $|C|$ is
a unit in $D$. The matrix
\[
                h=|C|^{-1}\sum_{t\in C}U(t)V(t)^{-1}
\]
reduces to the identity, hence is invertible. Permuting the sum by
left multiplication gives $U(t)h=hV(t)$ for each $t$. This
constructs one basis valid on the entire group, including at the
dyadic prime. In either case $\Psi=\psi_{\rm Gal}$ on this
inertia because $\varepsilon$ is unramified at $p$.

For an arbitrary place $w$ above $p$, choose a prolongation $B$
to $\Omega$. If $p\notin S$, the image of its absolute inertia
in $\Gamma$ is trivial, so one may use the three trivial
characters there in the level condition; the displayed $\xi_i$
also restrict trivially since $p\nmid N$. Otherwise $B$ is
conjugate over $\mathbb Q$ to $A_v$. Its inertia over $K$
embeds in the conjugate of $I^+$, because $G_K$ is the sign
kernel. Transport the preceding basis, together with the torus
conjugation from the integral model. For a positive-sign
conjugating element this is ordinary matrix conjugation. For a
negative-sign element the extra involution on a matrix $X$ is
$X\mapsto JX^{-t}J$, with $J$ reversing the coordinates. It
fixes $T(\Psi)$, and $\Psi$ is invariant under conjugation as a
character. Thus the conjugated inertia representation still has
the indicated diagonal in one basis on this same $B$. This
establishes every field of the native $\mathrm{LevelAway}$.
% END VUN2

% BEGIN VUN3
\paragraph{VUN3. The ordinary flag and its cyclotomic inequalities.}
\label{claim:VUN3}
Above $\ell$ the required weight vector is $(e,0,-e)$. On
$D_{v_\ell}\cap\ker(\eta\pi)$, the integral local basis makes
$r_D$ upper triangular with diagonal $\Psi^{d_i}$ on inertia.
Choose any $w\mid\ell$ of $K$ and any prolongation $B$ to $\Omega$.
Transport a local basis from the stored rational prolongation to $B$,
using conjugacy over $\mathbb Q$ as in VUN2. The group
$D_{K,B}$ identifies, inside $\Delta$, with
$D_{\mathbb Q,B}\cap\ker\eta_4$. Under a negative conjugation
$JX^{-t}J$ again preserves upper triangularity; it reverses and
inverts the three diagonal characters. On inertia this still
gives $T(\Psi)$, since this torus is fixed. Thus $\rho$ on
$D_{K,B}$ has one invariant complete ordered flag with characters
$\kappa_i$. On its inertia they satisfy
\[
                \kappa_i=\varepsilon^{e d_i}\xi_i
\]
with the same three global level-$N$ characters as in VUN2.

For the last disjunction in $\mathrm{LevelAt}$, we prove that
$\kappa_i\ne\kappa_j\varepsilon$ for every $i,j$, including
$e=1$. HSR3 gives a sign-positive element in the stored
decomposition group with $\varepsilon(g)=1$ and $\Psi(\pi g)^2-1$
a unit. Its conjugate at $B$ lies in $D_{K,B}$. At this element
the diagonal entries on the flag are units with pairwise distinct
residues: HSR4 proves this from the three distinct residual roots
of the integral characteristic polynomial. Replacing the three
roots by their inverses in reverse order preserves this property.
For $i\ne j$ they therefore differ at an element where
$\varepsilon=1$. For $i=j$, equality would force $\varepsilon$
trivial on $D_{K,B}$. This is impossible since the finite local
extension $K_w/\mathbb Q_\ell$ cannot contain all $\ell$-power
roots of unity, whose local degrees grow without bound.
This proves the unlinked alternative, independently of any
potential-crystallinity alternative.

Together VUN1--3 show that every $\rho$ supplied by the automorphy
witness of an eligible point satisfies the literal
$\mathrm{Adm3}(\iota,e,N,\rho)$ of RNC16. This specialization
uses the native point's local inertial basis and the host separator;
its local conclusions apply to both embeddings of the Gaussian field.
% END VUN3

% BEGIN VUN4
\paragraph{VUN4. One prime and one Frobenius element for every point.}
\label{claim:VUN4}
Take $m=\ell N$. Some prime divisor $p$ of $12m-1$ is $3$ modulo
$4$: the integer is $3$ modulo $4$, whereas a product of primes
$1$ modulo $4$ cannot have this residue. The prime $p$ divides
neither $3$ nor $m$. Hence $7\le p\le 12\ell N$ and $p\nmid
\ell N$. This prime is inert in $K\simeq\mathbb Q(i)$:
$X^2+1$ has no root in $\mathbb F_p$ but splits in $\mathbb F_{p^2}$.
The unique prime $w\mid p$ of $K$ has norm $q=p^2$, and the
discriminant of $K$ is $-4$, so $p\nmid\operatorname{disc} K$.

Fix a prolongation $B$ of the $p$-adic valuation to $\Omega$.
In the decomposition group over $\mathbb Q$ choose a lift of
arithmetic Frobenius and square it to obtain $\sigma$. Then
$\varepsilon(\sigma)=p^2$ and $\eta_4(\sigma)=1$.
For \emph{each} embedding $j:K\hookrightarrow\Omega$ the same
$\sigma$ fixes its image and defines $\tau_j\in G_K$. It belongs
to the decomposition group at $B$ and acts on residue elements
as the $p^2$-power map; it is an arithmetic Frobenius at $w$.
The two embeddings induce the same $w$ since $p$ is inert. Put
$\phi=\pi(\sigma^{-1})$. This one element of $\Gamma$ is fixed
before any representation or finite collection, even when
automorphy witnesses use different embeddings.

Set $z_0=\iota(\psi_{\rm Gal}(\sigma^{-1}))$, with
$O\hookrightarrow L$ understood, and $d=\varphi(N)$. The
Dirichlet character factors through $(\mathbb Z/N)^\times$, so
$z_0^d=1$. At the inverse Frobenius
$\Psi(\phi)=q^{-e}\psi_{\rm Gal}(\sigma^{-1})$. The three
diagonal terms of the original base point now give
\[
 \operatorname{tr}(r_z(\phi)^{\rm mat})
   = \iota^{-1}\!\left(q^{-e}z_0+1+q^e z_0^{-1}\right).
\]
This is precisely the geometric Frobenius convention of RNC16.
% END VUN4

% BEGIN VUN5
\paragraph{VUN5. The integer bound for one embedding of $K$.}
\label{claim:VUN5}
Fix the constants $a,b\in\mathbb N$ furnished by RNC16 for $K$,
independently of $\ell,c,N,e,\psi$ and the frame. At the prime
of VUN4 use the orbit $\mathcal O=\mu_d(\mathbb C)$, with each
root listed once and $d=\varphi(N)\ge1$. Every complex field
automorphism permutes this indexed orbit. The other catalogue
conditions hold: $p\ge7$, $p\nmid N\operatorname{disc}K$,
$\ell\ne p$ and $\ell$ is odd by $\mathrm{HostData}$.
Consequently RNC16 gives an integer $M\ne0$ with
\[
             \log |M|\le a\varphi(N)[N(e+1)]^b\log p .
\]
Its nonvanishing and product estimate hold at every subsequent
allowed coefficient realization. Here $p$ was selected using
$\ell$; this is an allowed specialization, since for this $p$
the catalogue works at every odd prime different from $p$.

Fix one embedding $j$ and a finite indexed family of eligible
points with automorphy witnessed through $j$. Assume that an
isomorphism of two members' matrices on $H$ implies equality
of their indices. Use the supplied restrictions $\rho_i$,
each satisfying the rank-three branch of the catalogue by VUN3.
Put their images into the one set $X_{z_0}$ and set all other
$X_z$ empty. By VUN1 an isomorphism of two restrictions would
identify the indices, so passing from the indexed family to its
image loses no factors. For each $i$ the value of $\rho_i$ at
$\tau_j^{-1}$ is exactly the matrix of $r_i$ at $\phi$.
By VUN4 the comparison trace in the catalogue is the original
base trace. Write $\delta_i=\operatorname{tr}(r_i(\phi)^{\rm mat})-
\operatorname{tr}(r_z(\phi)^{\rm mat})$. RNC16 gives all
$\delta_i\ne0$ and
\[
                |M|_\ell\le\prod_i|\delta_i|_\ell .
\]
The singleton case requires no prior pairwise hypothesis. In
particular the same $\phi$ separates \emph{every} allowed marked
type-I point satisfying $B'_s$.

Normalize $v_\ell(\ell)=1$. Since $M\ne0$ is an integer,
$\ell^{v_\ell(M)}$ divides $|M|$. Applying the negative
logarithm to the product bound therefore gives
\[
 \log\ell\sum_i v_\ell(\delta_i)
    \le -\log|M|_\ell
    \le \log |M|
    \le a\varphi(N)[N(e+1)]^b\log p.
\]
The product and sum conventions also give the stated inequality
for an empty fiber.
% END VUN5

% BEGIN VUN6
\paragraph{VUN6. The literal uniform host number.}
\label{claim:VUN6}
For the Gaussian field $K$, choose the exponent
$C=2a+b+4$ from VUN5 before all closures, primes, coefficient
rings, levels, characters, frames, complex embeddings and marks.
For any such data take $\phi$ from VUN4. For every later finite
list of marked characteristic-zero type-I points satisfying
$B'_s(K,\iota)$ and the precise pairwise nonisomorphism test
on $H$, choose one automorphy witness for each. Partition by the
associated field embedding $K\hookrightarrow\Omega$. There are
two embeddings, so summing VUN5 over at most two fibers gives
\[
             \log\ell\sum_i v_\ell(\delta_i)
               \le 2a\varphi(N)[N(e+1)]^b\log p .
\]
Each application uses just the one embedding of $K$ attached
to that fiber.

Put $X=\ell N e$. The host has $\ell\ge3$, $N$ positive and
even, and $e\ge1$, hence $X\ge6$. We have $2a\le X^{2a}$
(also for $a=0$), $\varphi(N)\le X$, $N(e+1)\le X$, and
$\log p\le p\le12\ell N\le X^3$. Multiplication proves
\[
  \left(\sum_i v_\ell\bigl(
     \operatorname{tr}(r_i(\phi)^{\rm mat})
      -\operatorname{tr}(r_z(\phi)^{\rm mat})\bigr)\right)\log\ell
       \le (\ell N e)^C .
\]
For all these nonzero arguments the source definition
$\mathrm{vl1}_\ell(x)=-\log|x|_\ell/\log\ell$ is precisely
$v_\ell(x)$. This inequality and the separate singleton trace
separation of VUN5 supply both clauses of the original
$\mathrm{HostNumber}_s$ for every marked $g_1$. The inherited
catalogue and automorphy comparisons are RNC16 and HSR8, used
on the native $\mathrm{Adm3}$ branch verified in VUN1--3.
% END VUN6

\section{Strong irreducibility and the marked rank-three automorphy implication}
\label{sec:markedaut-v64}

% BEGIN VMA1
\paragraph{VMA1. The absolute ordinary point on the prescribed host.}
\label{claim:VMA1}
Fix any $K$ with the source Gaussian CM condition $\mathrm{IsQi}$,
an algebraic closure $\Omega/\mathbb Q$, a prime $\ell$, the
comparison $\iota:\overline{\mathbb Q}_\ell\simeq\mathbb C$,
$c:\mathrm{Coef}(\ell)$, $N\ge1$, $e,\psi,o_1$ with
$\mathrm{HostData}$, and $S,(A_v),F$ with $\mathrm{IsHostFrame}$.
Put $O=c.O$, $k=O/\mathfrak m_O$, $E=\overline{\mathbb Q}_\ell$,
$\Delta=G_{\mathbb Q}$ on this $\Omega$, $\Gamma=G_S$ and
$p:\Delta\twoheadrightarrow\Gamma$. The pullbacks of the frame
characters are $\eta_4$ and $\Psi_0=\chi_\ell^e\psi_{\rm Gal}$.
Write $H=\ker\eta_4$, choose the frame's complex conjugation lift
$\sigma$ with $p\sigma=F.c$, and put $B=A_{v_\ell}$. Then
$\sigma^2=1$, $\eta_4(\sigma)=-1$, $\Psi_0(\sigma)=1$ by
complex conjugation and the given parity, and $v_B(\ell)<1$ in
multiplicative valuation notation. AFE10 supplies
$x\in D_B\cap H$ with $\Psi_0(x)^2-1\in O^\times$, including
at $\ell=3$. Thus a marked element exists already. In the
conditional joins below $g_*$ is the particular absolute marking
supplied jointly with the arithmetic ends, and $g_1=p(g_*)$.
In each case the identities $\eta_4(g_*)=1$ and
$\Psi_0(g_*)^2-1\in O^\times$ give $F.\mathrm{Marked}(g_1)$.

For every integral local coefficient algebra $A$ on which the point
conditions are defined, a marked point $r':\Gamma\to\mathrm{GG}_3(A)$
at $g_1$ gives the original absolute $\mathrm{OrdPt}$ at
$(\sigma,g_*,B)$ by pullback. The sign and residual congruences
pull back entrywise. The reductions are locally constant since
$p$ is continuous. The normalization at $\sigma,g_*$ follows from
$p\sigma=F.c$ and $p(\sigma g_*^{-1}\sigma)=F.c g_1^{-1}F.c$.
The prescribed upper triangular flag at $B$, including its inertial
diagonal $(\Psi_0,1,\Psi_0^{-1})$ on $H$, is the point's own
$\mathrm{CondT}$ after identifying the actual $D_B,I_B$.

For the away condition take any inertia element at a valuation
away from $\ell$ on which $\eta_4=\Psi_0=1$. The trivial
valuation has trivial inertia. Outside $S$ the quotient kills
the element. Inside $S$, conjugate the valuation to the chosen
$A_v$. At $v\in S_b$ the local diagonal condition makes all
diagonal entries one, hence the whole element one in its diagonal
basis. At the other $v\ne v_\ell$ the element is in
$I_v\cap\ker\eta_v\cap\ker\bar\Psi_v$, on which the point is
trivial. Conjugating back preserves the identity in $\mathrm{GG}_3$.
This verifies all five fields of the absolute ordinary point for the
same datum, in both residual and characteristic-zero specializations.
% END VMA1

% BEGIN VMA2
\paragraph{VMA2. Finite weight zero from the absolute finite end.}
\label{claim:VMA2}
Fix $g_*,g_1$ as in VMA1. Assume its $\mathrm{FinEnd}$: there is
one finite-index $V\le\Delta$ such that, over every complete DVR
$D$ with finite residue field and local $O$-structure annihilating
$\mathfrak m_O$, each absolute ordinary point on this datum has
the base characteristic polynomial on $V\cap H$. Let $R,u$ be
\emph{any} universal marked pair at $g_1$.
The host frame has $F.\mathrm{Ok}$ by the frame verifications in
CVI2: the decomposition separator is AFE10 for every $v\in T_b$,
away from $T_b$ the ramified part has order prime to $\ell$, and
the continuity, parity and Frobenius conditions follow from the
defining identifications of this host. These verifications use
$\mathrm{HostData}$ and $\mathrm{IsHostFrame}$ alone. The source
groups are profinite (the quotient kernel is closed), and $S$ is
finite. NMD3's marked-entry argument therefore gives a finite
adapted weighted presentation of this $R$ with variables
$z_i=u(h_i)_{a_i b_i}-(r_0(h_i))_{a_i b_i}$, $h_i\in\ker\eta$,
of weights $w_i=b_i-a_i$. They lie in $\mathfrak m_R$, and its
weight-zero component is $\mathrm{wpart}(O,u,0)$.

Take any local $O$-map $R\to D$ of the kind tested in UCE9.
It sends the universal point to an integral marked point:
congruences and local bases pass along the map, and
$\mathfrak m_R^n$ maps into $\mathfrak m_D^n$, giving continuity.
VMA1 gives its absolute ordinary realization. Since
$O\to D$ factors through $k$, the base character in the
$\mathrm{FinEnd}$ equality here is precisely the reduction
$\bar\Psi_0$. UCE10--11 apply with the original finite $k$,
$2\in k^\times$, the marked element $g_*$, the original normalization
and the given $V$. They yield one triangular direction on $H$ and
equality of all diagonal entries with the base. Explicitly, a positive
power of every element of $H$ lies in $V$; its eigenvalues are thus
algebraic over $k$. Adjoining an algebraic element of the local domain
gives a field over $k$, so a characteristic-polynomial coefficient
congruent to its base value in $k$ equals that base value. The marked
diagonal projectors and UCE11's trace-pairing argument then exclude
both a three-dimensional irreducible block and an irreducible outer
two-dimensional block. This works in characteristic three as well.

The map $p$ sends $H$ onto $\ker\eta$. The test map kills the $z_i$
of weight zero, and kills every $z_i$ on one of the two strict sides
of zero: off the diagonal the base entries vanish. Thus each
nonconstant monomial of weight zero vanishes at every such test
map. By the DVR-curve criterion UCE9 all these test maps imply
$\mathrm{wpart}(O,u,0)$ is a finite $O$-module. In particular
this conclusion is integral and allows nilpotents in $R$.
% END VMA2

% BEGIN VMA3
\paragraph{VMA3. The marked eigenvalue test on every open subgroup.}
\label{claim:VMA3}
Keep any mark $g_1$. Let $r:\Gamma\to\mathrm{GG}_3(E)$ be an
original characteristic-zero marked point of Type I. It has an
integral marked representative $r'$ over a finite DVR $A\subset E$
over $O$, related by conjugation by $T(t)=\operatorname{diag}(t,1,t^{-1})$
for one $t\in E^\times$. On $H$ the matrix of $r\circ p$ is an
ordinary representation $\rho$, as the sign there is positive.
It is irreducible because $p(H)=\ker\eta$ and Type I means
irreducibility on that kernel. The matrices of $r\circ p$ are
continuous: the local-constant reductions of $r'$ give continuity
in its finite coefficient field, and each matrix of $r$ is obtained
by fixed torus conjugation. For every open subgroup $U$ of
$\Delta$ contained in $H$, the restriction $\rho|_U$ is irreducible.

To prove this, take its finite-index normal core $U_0$ in $\Delta$.
For the chosen $B$ over $\ell$, take $s\in I_B$ with
$\chi_\ell(s)=1+\ell$. Indeed the cyclotomic extensions
$\mathbb Q_\ell(\mu_{\ell^n})$ are totally ramified, by the
Eisenstein polynomial for $\zeta_{\ell^n}-1$; the images on
inertia at each finite level are full, and compactness lifts the
compatible units to a single element. With
$m=[\Delta:U_0]\,|(\mathbb Z/N)^\times|$ we have
$u=s^m\in U_0\cap I_B$ and $\Psi_0(u)=(1+\ell)^{em}$.
The integer $em$ is positive. The three eigenvalues
$\Psi_0(u),1,\Psi_0(u)^{-1}$ in the ordinary basis are distinct.

Suppose $\rho|_{U_0}$ is reducible. A simple $U_0$-subspace exists
by minimal dimension. Its $H$-conjugates are simple and span the
whole irreducible $H$-space, so the restriction is semisimple.
The isotypic components form one $H$-orbit (an orbit sum is
$H$-stable). Their common dimension divides three. A reducible
restriction is therefore scalar, or consists of three distinct lines
permuted transitively: in the single isotypic case the multiplicity
times the dimension of the simple module is three. The scalar case
contradicts the three eigenvalues. The three lines in the second case
are the eigenlines of $u$. Its triangular matrix with distinct
diagonal entries can be diagonalized by an upper triangular basis
change, by choosing the eigenvector for the $i$th diagonal entry
in the $i$th flag step. Hence $H$ acts by monomial matrices in
this basis and the separator $x\in D_B\cap H$ of VMA1 remains
upper triangular there. Its permutation is the identity. Here the
separator is needed only in the decomposition group, and the
regular element $u$ is in inertia.

A transitive subgroup of $S_3$ contains a three-cycle, since its
order is divisible by three. Choose $g\in H$ mapping to this
cycle. Then $\operatorname{tr}\rho(gx^n)=0$ for $n=0,1,2$.
The equality also holds for the integral representative, by
conjugation on $H$. In its residue field put
$a=\overline{\Psi_0(g)}$, $b=\overline{\Psi_0(x)}$. Residual
diagonality yields $ab^n+1+(ab^n)^{-1}=0$ at these three indices.
Thus $a,ab,ab^2$ are roots of $X^2+X+1$. They are three distinct
elements since $b^2-1$ is a unit before reduction. This contradiction
proves irreducibility on $U_0$, hence on $U$.
% END VMA3

% BEGIN VMA4
\paragraph{VMA4. Irreducibility on the full cyclotomic kernel.}
\label{claim:VMA4}
For the same point, put $J=H\cap\ker\chi_\ell$. Then
$\rho|_J$ is irreducible over $E$. This subgroup can have infinite
index. Suppose the restriction is reducible. The sum of all
$H$-conjugates of a simple $J$-subspace is $H$-stable and is
all of $E^3$ by irreducibility. A finite direct sum of simple
subspaces can be extracted from this sum by successively adjoining
a simple subspace that increases the span. Thus $\rho|_J$ is
semisimple. It admits a nonzero proper equivariant projection $P$.
In $\operatorname{End}_J(E^3)$ the element
$P-(\operatorname{tr}P/3)I$ is traceless and nonzero.

The traceless part of this commutant is stable under conjugation
by $H$. Such conjugations commute on it, since $H/J$ is abelian.
A commuting family on a nonzero finite-dimensional space over the
algebraically closed field $E$ has a common eigenvector: if any
member is nonscalar, restrict to a nonzero proper eigenspace and
induct on dimension. Hence a nonzero traceless matrix $M$ satisfies
\[
       \rho(h)M\rho(h)^{-1}=\varphi(h)M\qquad (h\in H)
\]
for a character $\varphi:H\to E^\times$. The kernel of $M$ is
$H$-stable, so $M$ is invertible. Taking determinants gives
$\varphi^3=1$. Its kernel is the commuting locus of $M$ and the
continuous $\rho(h)$, hence is closed. A closed finite-index
subgroup is open (its complement is a finite union of closed
cosets), and $H$ itself is open in $\Delta$. By VMA3 the kernel
acts irreducibly. The eigenspace of any eigenvalue of $M$ is stable
under that kernel and nonzero, hence all of $E^3$. This makes
$M$ scalar, contrary to its nonzero trace-zero condition in
characteristic zero. The argument also applies when $\ell=3$.
% END VMA4

% BEGIN VMA5
\paragraph{VMA5. The reduced and full marked assertions are equivalent.}
\label{claim:VMA5}
The reduced source assertion $\mathrm{ThmS6Core}_s$ supplies,
after the host, frame and $\iota$ are fixed, one marked $g_1$.
For every universal $R,u$ at this mark it gives finite weight zero;
for every $O$-algebra map $\xi:S(u)\to E$ from the graded
entry subalgebra and every group map $r:\Gamma\to\mathrm{GG}_3(E)$
satisfying $\mathrm{IsPtHom}(u,\xi,r)$ and Type I it gives
$\mathrm{AutOnAlt}(K,\iota,-1,(r\circ p)^{\rm mat})$.
With exactly these quantifiers it implies $\mathrm{ThmS6}_s$.

Indeed keep that mark and the supplied finiteness for the given pair.
The characteristic-zero realization VSCM8 applies to this universal
ring with its adapted presentation from NMD3. Its hypotheses are
verified on this instance by the profinite quotient, $F.\mathrm{Ok}$,
the mark and the supplied weight-zero finiteness. Concretely the
weight-zero coordinates land integrally, the nonconstant mixed-sign
monomials of weight zero land in the maximal ideal, and one scalar
$t\in E^\times$ places all the rescaled marked coordinate differences in the maximal ideal
of a single finite coefficient DVR $A$. Evaluation of the homogeneous
relations extends to the completed local ring. The resulting
integral $r_A$ satisfies
\[
     \mathrm{push}_{A\hookrightarrow E} r_A
                    =\mathrm{conj}_{T(t)} r.
\]
The entry and sign identities in $\mathrm{IsPtHom}$ give this equality
for the particular $r$ in the quantifier. This is
$F.\mathrm{IsCharZeroPt}(g_1,r)$. VMA3--4 now provide the two
additional fields $\mathrm{opens}$ and $\mathrm{cyc}$ of $B'_s$ on
the actual absolute Galois subgroups. Its third field is precisely
the supplied $\mathrm{AutOnAlt}$ for this same $\xi,r$ and comparison
$\iota$. This finishes the full implication. Conversely, retaining the
given mark and finite module and projecting $B'_s$ to its third field
gives $\mathrm{ThmS6Core}_s$ immediately.
% END VMA5

% BEGIN VMA6
\paragraph{VMA6. An arithmetic end with one mark suffices.}
\label{claim:VMA6}
The original absolute $\mathrm{HostEndAlt}$ implies
$\mathrm{ThmS6}_s$. To see this, choose the specified $\sigma,B$ of
VMA1 after fixing the original host data, its frame and comparison.
The arithmetic end now supplies \emph{one} absolute $g_*$ together
with $\mathrm{Marks}$, $\mathrm{FinEnd}$ and $\mathrm{AutEnd3Alt}$
at $(\sigma,g_*,B,K,\iota)$. Fix $g_1=p(g_*)$ before $R,u$,
the coordinate map or any representation. VMA1 descends the
marking. For each universal $R,u$ its weight-zero module is finite
by VMA2. For each subsequent coordinate point $r$ of Type I, take
the single integral representative $r_A$ just constructed in VMA5.
The proof of this realization uses only finiteness, independently
of any automorphy assertion. VMA1 applied to $r_A$ gives exactly
the $\mathrm{OrdPt}$ needed by $\mathrm{AutEnd3Alt}$. On every open
$U\subset H$, its matrices over $E$ are conjugate by $T(t)$ to
those of $r\circ p$, so VMA3 gives irreducibility there.

Choose an embedding of $K$ in $\Omega$ over $\mathbb Q$. Its
image is $\mathbb Q(i)$ and the image of its absolute Galois group
in $\Delta$ is $H$, since the action on $i$ defines $\eta_4$.
Let $\rho$ and $\rho_A$ be the resulting ordinary
$\mathrm{GL}_3(E)$ representations of $G_K$ obtained from
$r\circ p$ and $\mathrm{push}(r_A)\circ p$. They have literally
the specified matrices on restriction; their signs there are
positive. The conjugation $T(t)$ intertwines them, and still
intertwines their twists by $\chi_{K,\ell}^{-1}$.
Apply $\mathrm{AutEnd3Alt}$ to the absolute integral point and
$\rho_A$ to get $\mathrm{IsAutCsdA}$ for that twisted representation.
Transport this \emph{full} assertion to the twist of $\rho$.
For the first conjunct, the same regular algebraic
conjugate-self-dual cusp works: conjugation carries invariant
complements and preserves every local Frobenius polynomial of the
attachment. For the second conjunct, continuity is invariant under a
fixed conjugation; carrying the local flags retains their characters,
hence both $\mathrm{Unl1}$ and $\mathrm{PotUnl}$. For every cusp
satisfying the stated attachment, use the same real $t_0$ and the
same places in its $\mathrm{WinAt}$ assertion. This gives the
complete conditional $\mathrm{WitA}$ of the original Alt conclusion.
The chosen $K$-embedding, its scalar tower, $\rho$, and its literal
matrix identification therefore witness $\mathrm{AutOnAlt}(K,\iota,-1)$
for $r\circ p$. Apply VMA5 with this finite module and automorphy
at the unchanged mark to obtain $\mathrm{ThmS6}_s$.

What remains to obtain the whole selected assertion by this route is
the simultaneous supply of $\mathrm{FinEnd}$ and $\mathrm{AutEnd3Alt}$
at that mark for every original host. In $\mathrm{FinEnd}$ the
finite-index subgroup precedes all residual DVR points. In
$\mathrm{AutEnd3Alt}$ the conclusion is the inverse-cyclotomic
$\mathrm{IsAutCsdA}$ for every integral ordinary point satisfying the
open-subgroup hypothesis, at the fixed $\iota$. AFE16 and AFE21
supply an arithmetic frame and the implication from its rank-three
$\mathrm{FinPt}$ to $\mathrm{FinEnd}$ at their given cohomological
input. The pinned arithmetic route to the joint end also calls upon
the linked automorphic end, the rank-three local comparison at even
level, and conjugate-self-dual descent. VMA1--6 establish the
comparison from the joint end to the literal selected assertion;
the joint arithmetic end remains an input.
% END VMA6

\section{A tame tower for the prescribed Case III label}
\label{sec:tensor-tower-case-iii}

The Case III argument needs four distinct residual diagonal characters
at one element of the field on which the argument will take place. It
also needs a quadratic field that realizes one particular ratio of
these characters. We construct both fields and the element together.

% BEGIN TFT1
\paragraph{TFT1. Two character values available before changing fields.}
Fix the original algebraic closure $\Omega/\mathbb Q$, prime $\ell$,
coefficient datum $c$, and raw datum $D$ satisfying $\mathrm{Ctx3}$.
Write $O=c.O$, $k=O/\mathfrak m$, $\Delta=G_{\mathbb Q}$ inside
this closure, and use bars for reduction to $k$. The finite character
$\psi$ here is the Galois character of the host Dirichlet character.
Let $m$ be the order of the latter. AX1 and the Case III conditions give
\[
 \ell\ge5,\quad \ell-1\mid 2e,\quad
 m>1\text{ odd},\quad (m,\ell)=1,\quad \psi^m=1.
\]
Indeed $\psi$ is nontrivial by Case III and $m$ divides the odd
part of the host order $o$, which is prime to $\ell$.
Thus no inertia image of $\psi$ has order divisible by $\ell$,
so $D.S_b=\varnothing$. The conductor of $\psi$ divides the level
$N$, prime to $\ell$. The character $\xi$ is the product of an odd
dyadic character of modulus $2^n$ and an auxiliary character of
modulus $p_1$, with $n=n_{\rm III}$, $p_1\ne\ell$ an odd
prime coprime to $N$. Thus $\xi$ and $\psi$ are trivial on
every inertia group above $\ell$. Write $\omega$ for the
Teichm\"uller lift of the residual cyclotomic character.
Reduction is injective on roots of unity of order prime to $\ell$
since $1+\mathfrak m$ is pro-$\ell$.

There are $a\in I_2$ at the specified dyadic place and $u\in\Delta$
with the following values, where $d$ has nontrivial odd order dividing $m$:
\[
 \begin{array}{c|ccc}
      &\xi&\psi&\omega\\ \hline
 a&-1&1&1\\
 u&1&d&1
 \end{array}
\]
At $2$, inertia surjects onto the units modulo $2^n$ because the
dyadic cyclotomic field is totally ramified there. Take an element
acting by $-1$ on $\mu_{2^n}$. It is trivial on the odd conductor
of the auxiliary factor of $\xi$ and on $\mu_\ell$, while the dyadic
factor is odd. Raising to the odd power $m$ gives $a$.
For $u$, pick a unit modulo $N$ on which $\psi$ is nontrivial.
Since $(N,\ell p_1)=1$, the Chinese remainder theorem and the
cyclotomic Galois action give a lift acting by that unit on $\mu_N$
and trivially on $\mu_{\ell p_1}$. Its $2^n$-th power kills the
dyadic factor of $\xi$ and preserves a nontrivial $\psi$-value,
as $m$ is odd. This is $u$.
% END TFT1

% BEGIN TFT2
\paragraph{TFT2. The real field and two distinct quadratic roles.}
Let $S=D.S$, which contains $2$ and $\ell$. Choose a new prime
$q>\ell$ larger than all prime divisors of $2N\ell p_1$ and than
all members of $S$, with $q\equiv-1\pmod{8\prod_{v\in S}v}$.
Dirichlet's theorem applies to this reduced residue class. The
quadratic field $K'=\mathbb Q(\sqrt{-q})\subset\Omega$ has
discriminant $-q$ and associated character $\eta$. Each odd
$v\in S$ splits since $-q\equiv1\pmod v$; $2$ splits since
$-q\equiv1\pmod8$. Inertia at $q$ contains an element $b$
with $\eta(b)=-1$. At this prime $\xi,\psi,\omega$ are all unramified.
The sign character $\eta$ is regarded also as an $O$-valued unit.

Let $\Psi=\chi_\ell^e\psi$ and form the open subgroup
\[
 H=\{g\in\Delta:\ \xi(g)=\eta(g),\quad
                          \bar\Psi(g)=\bar\psi(g)\}.
\]
The two character quotients have finite image and open kernel;
the second is $\bar\omega^e$. Thus $H$ is normal with abelian quotient.
All complex conjugations lie in it: $\xi$ and $\eta$ are odd,
$\Psi$ is even by goodness, and $\psi$ of odd order kills
complex conjugation. For $t$ in any inertia group over $\ell$,
$t^2\in H$. Indeed $\xi(t)=\psi(t)=1$, $\eta(t)^2=1$ and
$\bar\omega(t)^{2e}=1$ since $\ell-1\mid 2e$.

Choose two further, distinct quadratic characters $\lambda,\lambda'$
of negative prime discriminants, by the same progression construction.
For $\lambda$ impose splitting at $S\cup\{q\}$; for $\lambda'$
impose splitting there and at the ramification prime of $\lambda$.
Let $K$ be the field of $\lambda$, put
$J=\ker\lambda\cap\ker\lambda'$, and let $F$ be fixed by
$H\cap\ker(\lambda/\lambda')$. Then
\[
 G_F\subseteq H,\qquad G_F\cap G_K=H\cap J.
\]
Here and below we identify these Galois groups with subgroups
of $\Delta$. The field $F$ is finite abelian and totally real since
both auxiliary characters are odd and $H$ contains complex conjugation.
Its degree is even: $\lambda/\lambda'$ is nontrivial on inertia at the
prime of $\lambda$, and its quadratic field lies in $F$.
The field $K$ is imaginary quadratic and splits every prime of $S$.
Put $L=FK$. At the prime of $\lambda$, inertia inside $G_F$
has $\lambda=\lambda'=1$; away from that prime $\lambda$ is unramified.
Consequently $L/F$ is unramified at every finite place.
Every square in $\Delta$ lies in $J$, and $J$ contains every absolute
decomposition group over $S\cup\{q\}$. In particular
$t^2\in G_L$ for all inertia elements $t$ over $\ell$.
% END TFT2

% BEGIN TFT3
\paragraph{TFT3. One separating element for either allowed label.}
Fix the supplied Boolean label $\mathrm{pr}$ satisfying
$\mathrm{LabOk}(D,\mathrm{pr})$. Set $\epsilon=1$ for the false
label and $-1$ for the true label. Both $a$ and $b$ of TFT1--2
belong to $J$ by the splitting at $2,q$, and $u^2\in J$.
Put $g_0=a b u^2$. Then
\[
 \xi(g_0)=\eta(g_0)=-1,\qquad
 \omega(g_0)=1,\qquad \psi(g_0)=d^2 .
\]
In fact $\eta(a)=1$ by splitting at $2$, and $\xi,\omega,\psi$
are trivial on $b$. Thus $g_0\in H\cap J=G_L$. Write
$z=\bar d^{\,2\epsilon}$. This is of nontrivial odd order, since
reduction preserves the order of $d$ and its order is odd.
In particular $z\ne 1,-1$.

For the relabelled four diagonal characters $\chi_i$, their quotients
$\chi_0/\chi_i$ in order are $1,\alpha,\Psi^\epsilon,
\alpha\Psi^\epsilon$, where $\alpha=\chi_\ell\omega^{-1}\xi$.
This follows directly from the two permutations of the diagonal.
On $G_F$, $\bar\alpha=\bar\xi$ and $\bar\Psi=\bar\psi$.
Consequently the residual root at $w=(h,k')$ is
\[
 r_w=\bar\xi^h(\bar\psi^\epsilon)^{k'},\qquad
 r_w(g_0)=(-1)^h z^{k'}.
\]
At $g_0$ the reduced four quotients are $1,-1,z,-z$, all distinct
as $\ell\ge5$. Thus the four diagonal values themselves are distinct.
The eight roots are $\pm(1,0),\pm(0,1),\pm(1,1),\pm(-1,1)$.
Their values are among $-1,z^{\pm1},-z^{\pm1}$ and are all
nontrivial. The root $(1,1)$ has value $-z$ of order
$2\operatorname{ord}(z)\ge6$, so no power from one through four kills it.

Suppose a quadratic character $\nu$ of $G_L$ is trivial on
$G_L\cap\ker\omega$ and multiplication by $\nu$ permutes the four
restricted diagonal characters in the sense of the source. Apply this
condition to $\bar\chi_0$. The resulting equality
$\nu\bar\chi_0=\bar\chi_j$ at $g_0$ has $\nu(g_0)=1$, so
separation gives $j=0$. The same equality on all of $G_L$ gives
$\nu=1$. All three character tests use the same field and the
same $g_0$ for the initially prescribed label.
% END TFT3

% BEGIN TFT4
\paragraph{TFT4. Inertia and the local alternatives on this tower.}
For every place of $\Omega$ over $\ell$, the mod-$\ell$ cyclotomic
image on rational inertia is $\mathbb F_\ell^\times$, since
$\mathbb Q_\ell(\mu_\ell)/\mathbb Q_\ell$ is totally ramified of
degree $\ell-1$. Choose an inertia element $t$ mapping to a generator.
Its square has $\bar\omega(t^2)\ne1$ because $\ell-1\ge4$,
and $t^2\in G_L$ by TFT2. The relative inertia is the intersection
of absolute inertia with $G_F$, and $\xi(t)=\psi(t)=1$. Hence
$r_w(t^2)=1\ne\bar\omega(t^2)$ for \emph{every} $w\in\mathbb Z^2$.
This gives both the cyclotomic exclusion on $G_L$ and, at every
place over $\ell$, the first alternative of $\mathrm{locL}$ on
relative inertia. At those places all inertia elements, on squaring,
fix $F$; the uniform exponent $2$ is prime to $\ell$ and gives
exactly $\mathrm{TameL}(\ell,F)$.

The original diagonal residual representation $\bar r_0$ is trivial
at an element precisely when $\bar\alpha=\bar\Psi=1$: its diagonal
characters are $\alpha\Psi,\Psi,\alpha,1$ before relabelling.
On relative inertia over $\ell$, $\bar\alpha=\bar\xi=1$ and
$\bar\Psi=\bar\psi=1$ by TFT2. Away from $\ell$,
$\omega$ is unramified, hence $\bar r_0=1$ on an inertia element
implies $1=\bar\Psi=\bar\psi$ there. Since $\psi^m=1$ and
$(m,\ell)=1$, this gives $\psi=1$ in $O$. Finally, $\psi^m=1$
holds on all of $G_F$, in particular on its full cyclotomic kernel.
These are the integral and residual inertia clauses of $\mathrm{Ok}$.
% END TFT4

% BEGIN TFT5
\paragraph{TFT5. The original tensor-factor tower exists.}
Use $F,K,L$ from TFT2 and $g_0$ from TFT3. The compositum $L=FK$
is CM, quadratic over its totally real subfield $F$, since $K$
is imaginary quadratic. For every prime of $F$ above $\ell$ take
one prolongation to $\Omega$ whose restriction to $K$ is that of
the prescribed valuation $D.\mathrm{pl}(\ell)$. Such a choice is
possible because the two primes of $K$ over $\ell$ split after
base change to each completion $F_v$, and the involution of $L/F$
interchanges the two local factors. Absolute Galois over $F$ acts
transitively on the prolongations of each $v$. These finitely many
representatives therefore give exactly the original orbit and
place-on-$K$ requirements. Together with the supplied label these
data define the tower $T$ inside the original $\Omega$.

For $T.\mathrm{Ok}$, the label hypothesis is the one supplied, and
TFT2 gives even degree, unramified $L/F$ and splitting of $S$ in
$K$. TFT3 gives separation at $T.g_0$, nontriviality of each
root, the order test and exclusion of the quadratic self-twists.
TFT4 gives the remaining root, inertia and local tests.
For $\mathrm{Four}(T)$, TFT1 gives $\ell\ge5$, the empty bad set
and a prime-to-$\ell$ exponent for the original $\psi$.
Its quadratic witness is $K'=\mathbb Q(\sqrt{-q})$ of TFT2.
It is totally complex of degree two and $\ell\nmid\operatorname{disc}(K')$.
The decomposition groups above $\ell$ and all primes of $D.S$
fix $K'$. For $g\in G_F$ the root $w_\alpha=(1,0)$ has value
$\bar\xi(g)=\bar\eta(g)$, hence equals $1$ on $G_F\cap G_{K'}$
and $-1$ on its complement. This supplies the two labelled
restrictions of the original $\mathrm{QuadOn}$ record.
The same tower is tame by TFT4. We have proved, for every original
$D$ in Case III and every permitted $\mathrm{pr}$, the existence of
\[
 T\quad\text{with}\quad T.\mathrm{pr}=\mathrm{pr},\quad
 T.\mathrm{Ok},\quad \mathrm{Four}(T),\quad
 \mathrm{TameL}(\ell,T.F).
\]
This supplies the initial tower used in CJP3. Producing the seeded
pair on a later tower belongs to the separate rank-four reduction.
% END TFT5

\section{Four polynomial links between the lattice and the Case III seed}

Fix an original Case III datum $D$, an acceptable, tame tower $T_0$
with the source four-factor property, an embedding $\iota$, the ambient
residue field $k$, and a global residual pair $c_0$ in the prescribed
lattice shape. Assume the cohomological input, the fitness implication
and the seed hypothesis $\mathrm{SeedIVAlt}$. Our aim is one extension
$T_1$ of $T_0$ with a property for \emph{every} later deep tower $T$:
on $T$ one integral seed must work for every admissible restriction of
$c_0$. The seed hypothesis gives a residual pair before $T_1$ is
chosen. The integral seed will only be chosen on $T$.

Use zero-based matrix indices. Write $\alpha,\beta,\sigma,\tau$ for
the source weights $\mathrm{wAl},\mathrm{wBe},\mathrm{wPs},\mathrm{wTh}$.
The support sets for the two types of lines and their intersection are
\[
 N_P=\{\beta,-\alpha,-\tau,-\sigma\},\qquad
 N_M=\{-\alpha,-\beta,-\sigma,-\tau\},\qquad
 Q_0=\{-\alpha,-\sigma,-\tau\}.
\]
Support restricts the off-diagonal entries of the representation.
Saying a pair is in $Q_0$ also restricts its local lower unitriangular
flag matrices. Write $E_\lambda$ for the existence of a nonzero
representation entry of weight $\lambda$. Two pairs have
$\mathrm{SamePs}$ if the $20$ entries agree both on the group and in
each local flag. $\mathrm{SameAl}$ means the corresponding equality of
$10$ entries.

% BEGIN RFCR1
\paragraph{RFCR1. Preserving one simple root while introducing the other.}
On the fit initial tower apply the seed hypothesis at the given
$\iota,k$ (with the original ambient datum). It supplies a split prime
$w$ and a residual pair $s$ supported on $N_M$ with auxiliary set
$\{w\}$. We construct $p_2$ on a larger finite split set $W_1$, in
$Q_0$, with $\mathrm{SameAl}(p_2,s)$ and $E_{-\sigma}(p_2)$.

Put $x(g)=s.r(g)_{10}/\bar\chi_0(g)$ and $m_v=s.n_{v,10}$, where
$\bar\chi_i$ are the actual residual diagonal characters. Lower
triangular multiplication makes $x$ a cocycle for
$\rho=\bar\chi_1/\bar\chi_0$. Normalization gives $x(g_0)=0$.
The first column of the local ordinary flag gives
$x(h)=(1-\rho(h))m_v$ on the decomposition group. We need an
independent cocycle $y$ for $\nu=\bar\chi_2/\bar\chi_0$.

Here is the interchange of roots in the cup construction of UX3.
Let $k_0$ be the image of the finite coefficient residue field, $L$ the
finite $k_0$-span of the values of $x$. This span is finite-dimensional
by compactness and local constancy. On $M=k_0\oplus L$ let
\[
 g(a,b)=\bigl(\nu(g)a,\ \rho(g)\nu(g)b+\nu(g)a\,x(g)\bigr).
\]
The cocycle rule for $x$ gives the action law. Use the finite-module
argument (UX.15) for $M$ and $(1,0)$. At the old places impose the
zero local class. At each added splitting prime allow the additional
one-dimensional tame inertia image spanned by $(1,0)$. Adding enough
primes makes the lower bound for the Selmer cardinality exceed the
number of cocycles unramified at all the added primes. The resulting
cocycle $(y,z)$ takes a nonzero multiple of $(1,0)$ on some such
inertia, where both characters are trivial. Subtract the coboundary
\[
 g\longmapsto\bigl((1-\nu(g))b,\
 (1-\rho(g)\nu(g))c-\nu(g)b\,x(g)\bigr)
\]
with $b=y(g_0)/(1-\nu(g_0))$ and
$c=z(g_0)/(1-\rho(g_0)\nu(g_0))$. Standardness separates both
roots at $g_0$. This sets both coordinates to zero there and preserves
the witnessed nonzero $y$. Write $n_v$ for its resulting local
constant, so $y(h)=(1-\nu(h))n_v$.

The cross term of $z(gh)$ is $x(g)\nu(g)y(h)$. To put the coordinates
back in the original matrix order set $z'=xy-z$. Expanding
$x(gh)y(gh)$ gives cross term $y(g)\rho(g)x(h)$ for $z'$.
The same subtraction at a local group replaces its central constant
$o_v$ by $o'_v=m_v n_v-o_v$. Use the matrices
\[
 U(x,y,t)=\begin{pmatrix}1&0&0&0\\ x&1&0&0\\ y&0&1&0\\ t&y&-x&1\end{pmatrix},
 \qquad r(g)=U(x,y,2z'-xy)\operatorname{diag}(\bar\chi_i(g))
\]
and the local matrices $U(m_v,n_v,2o'_v-m_vn_v)$. Their product
adds $yx'-xy'$ to the two central entries, exactly the cup identity
just obtained. The local identity gives the prescribed diagonal
characters after conjugation. The matrices are symplectic up to the
given diagonal multiplier, normalized at $g_0$, and trivial on the
required forbidden inertia. They define the promised $p_2$: the
$10$ data are unchanged and $r_{20}=y\bar\chi_0$ is nonzero somewhere.
% END RFCR1

% BEGIN RFCR2
\paragraph{RFCR2. Two simultaneous modifications before choosing the later tower.}
First restrict $c_0$ to $T_0$ and move the global ordinary flag to
the chosen places over $\ell$. More explicitly, if $g$ moves the
chosen place to a place above the new one, factor $c_0.r(g)c_0.n$
into a lower unitriangular matrix and an upper matrix. The successive
leading minors are units by the lattice support, as in UX2 and UX11.
Use the lower factor as flag; the upper factor conjugates upper matrices
to upper matrices and preserves their diagonal characters. Restriction
preserves the multiplier and continuity. The chosen $g_0$ is the
initial separating element. The ramification, inertia and diagonal
conditions of the restricted frame follow from those of the lattice
and fitness. This supplies $p_0$ over $c_0$, residual on the empty
auxiliary set, with representation supported on $N_P$ and flags on
$Q_0$. In particular the lower factor has zero $21$ entry (UX.13).

Regard $p_0,p_2$ on $W_1$ by weakening the forbidden-inertia test.
Each is supported on $N_P$ with flags on $Q_0$. For both take the
$20$ entry divided by $\bar\chi_0$ and its $20$ flag coordinate.
These are normalized $-\sigma$ cocycles and their local constants
by the first-column identity. Apply the \emph{simultaneous}
construction (UX.14)--(UX.19) to the two coordinates. It produces
one $W_2\supset W_1$, one nonzero normalized $-\alpha$ cocycle
and its local constants, and two $Q_0$ pairs $p_3,p_4$. Thus
\[
 \mathrm{SamePs}(p_0,p_3),\qquad
 \mathrm{SameAl}(p_3,p_4),\qquad
 \mathrm{SamePs}(p_4,p_2),\qquad E_{-\alpha}(p_3).
\]
In particular equality of the middle $10$ entries includes their local
constants. Choosing the two modifications independently would not
give this equality. All primes of $W_2$ split in the precise sense
of $T_0$.

Choose $T_1$ now. For each $w'\in W_2$ adjoin to the totally real
tower field the abelian degree-$\ell$ characters unramified outside
$w'$ by taking the compositum of their fixed fields of kernels.
Here $w'\ne\ell$. At $\ell$ the rational inertia group acts
nontrivially on $\mu_\ell$, whereas splitting requires the
Teichm\"uller cyclotomic character to be trivial there.
The character fields lie in the cyclic degree-at-most-$\ell$ part
of $\mathbb Q(\mu_{w'})$ by reciprocity. Their compositum over the
initial field is totally real, of exponent dividing $\ell$ and
unramified over $\ell$. Use the same label and auxiliary base field
and the $\ell$th power of the marked element in the new tower
(the tower extension construction in UX1). These operations require
acceptability and tameness of the initial tower and the specified
splitting primes, all present here. The tower remains tame at $\ell$;
its field kills every permitted degree-$\ell$ character at every
$w'\in W_2$. This is the source condition $\mathrm{Dies}$, also
valid at all later towers.
% END RFCR2

% BEGIN RFCR3
\paragraph{RFCR3. Restriction preserves the two matched coordinates and their witnesses.}
Fix $T\ge T_1$. The tower relation is transitive and preserves
fitness and acceptability. Restrict each $p_i$ for $i=0,2,3,4$ to
the new group and use the lower factors of its moved flags as in
RFCR2. Call the resulting pairs $q_i$. For the $N_M$ shape the
matrix product of a lower triangular representation and the old
lower flag is already lower triangular with invertible diagonal.
For the $N_P$ shape the conjugated flag has its lower factor in
$Q_0$ by the same minor calculation. Changes of the element moving
a place only multiply the moved flag on the right by an upper
matrix, so its normalized lower factor is unique. The original
ambient residual field and the four characters stay fixed.

At each auxiliary prime all diagonal characters on inertia are $1$.
The two allowed shapes are triangular in their respective orderings,
so $(r(g)-1)^4=0$ on inertia. Here $\ell\ge5$ by the four-factor
hypothesis, hence $r(g)^\ell=1$. The image is finite by local
constancy and compactness. It is an $\ell$-group (each of its elements
has order dividing $\ell$), hence kills wild inertia and factors
through the cyclic order-$\ell$ quotient of tame inertia. Splitting
gives $w'\equiv1\bmod\ell$ and identifies the local base field with
$\mathbb Q_{w'}$; the cyclotomic degree-$\ell$ character at $w'$
is nontrivial on this quotient. The $\mathrm{Dies}$ condition thus
kills the residual inertia on $T$. Away from $W_2$ the old forbidden
inertia was already killed. The new marked element restricts to an
$\ell$-power of the old one, so is still diagonal; the new local
characters are those at the original places by conjugacy. This gives
the residual conditions with the empty auxiliary set.

For an $N_P$ pair the first column of the factorization gives
\[
 n'_{20}r_{00}(g)=r_{20}(g)+r_{22}(g)n_{20}.
\]
For an $N_M$ pair one similarly has
$n'_{10}r_{00}(g)=r_{10}(g)+r_{11}(g)n_{10}$.
Both identities use the actual place mover. Since the diagonal entries
are the common nonzero characters, the first identity transports
$\mathrm{SamePs}$ and the second transports $\mathrm{SameAl}$.
A second restriction of the same global pair also has exactly the
same lower flags, by uniqueness of those factors. Consequently any
residual $u$ above $c_0$ with the stipulated flag compatibility is $q_0$.

We record why the needed nonzero entries survive. On the indicated
supports, the entry $d=r_{20}$ for $N_P$, or $d=r_{10}$ for $N_M$,
satisfies
$d(gh)=a(g)d(h)+d(g)b(h)$ with its two diagonal characters $a,b$.
Take $x$ with $d(x)\ne0$, and $y$ with $a(y)\ne b(y)$ by separation.
If $a(x)=b(x)$ and $d(y)=0$, use $z=yx$; if $d(y)\ne0$ in this
case use $z=y$. Otherwise take $z=x$. In each case $d(z)\ne0$ and
$a(z)\ne b(z)$. For $N=\ell^m$ the power recurrence gives
\[
 d(z^N)\bigl(a(z)-b(z)\bigr)
       =d(z)\bigl(a(z)^N-b(z)^N\bigr)
       =d(z)\bigl(a(z)-b(z)\bigr)^N\ne0 .
\]
The tower extension contains the $\ell^m$th powers of the old Galois
group in the image of restriction for some $m$. This proves
$E_{-\alpha}(q_3)$ and $E_{-\sigma}(q_2)$ from the particular witnesses
constructed in RFCR1--2.
% END RFCR3

% BEGIN RFCR4
\paragraph{RFCR4. The seeded four-link chain on every later deep tower.}
Take any later deep $T$. Apply the seed promise of RFCR1 at this $T$,
with the original prime $w$: it lies in $W_2$ and hence dies on $T$.
This gives a finite coefficient extension and one integral seed pair
at $\iota$. Write $z$ for its reduction into the original $k$. Its
residual pair is obtained by reducing the integral pair and composing
the finite residue-field embedding into $k$. The coefficients agree
by compatibility of the original ambient maps; reduction mod the
maximal ideal gives local constancy. Thus $z$ satisfies the residual
conditions on the same empty frame. It lies above $s$ and satisfies
$\mathrm{InA}$, $E_{-\beta}$ and $E_{-\tau}$. Here
$\mathrm{InA}$ means representation support
$\{-\beta,-\sigma,-\tau\}$ and a nonzero $10$ coordinate of every
local flag. This same $z$ works for all the permitted residual $u$.

Take such a $u$. The initial global lattice shape and the prescribed
flag equations give either
\[
 (E_{-\sigma}(u),E_\beta(u))\quad\text{or}\quad
 (\mathrm{InV}(u), E_{-\alpha}(u)).
\]
In the second alternative $\mathrm{InV}$ means representation support
$\{\beta,-\alpha,-\tau\}$ with flags supported on $Q_0$ and all
$n_{20}\ne0$. To recall the restriction argument of UX2, a normalized
single-root cocycle that vanishes on the tower subgroup also vanishes
on conjugates of its separating element. The cocycle identity then
forces it to vanish globally, since the two diagonal values there
differ. In the second lattice alternative the flag equation in (UX.13)
propagates $n_{20}\ne0$. If the restricted $10$ entry vanished
identically, the global $10$ entry would vanish by the same argument;
the other first-column identity would make $n_{10}/n_{20}$ solve the
local equation excluded in the global lattice shape. Indeed on the initial
decomposition group at $\ell$, the $20$ flag equation reads
$n_{20}\bar\chi_0=r_{20}+\bar\chi_2 n_{20}$. Here $r_{20}=0$
and $n_{20}\ne0$, so $\bar\chi_0=\bar\chi_2$. This proves the
stated alternatives for the original $u$.

By RFCR3 this $u$ is $q_0$. The successive matching conditions are
\[
 u\ \xrightarrow{\mathrm{SamePs}}\ q_3\
   \xrightarrow{\mathrm{SameAl}}\ q_4\
   \xrightarrow{\mathrm{SamePs}}\ q_2\
   \xrightarrow{\mathrm{SameAl}}\ z .
\]
The last equality uses that $q_2$ restricts $p_2$, $z$ lies above $s$,
and RFCR1 preserves the $10$ data of $s$. We check goodness of every
endpoint. For $N_P$ support, $E_{-\sigma}$ is detected by the $20$
entry: the multiplier equation gives zero in $31$ whenever $20=0$.
For $N_M$ support the corresponding test for $E_{-\alpha}$ is the
$10$ entry, using $32$ as its reflected partner.
\begin{itemize}
\item At $u$ the alternatives just proved each imply ordinary goodness.
At $q_3$ the $-\alpha$ witness survives by RFCR3. If $u$ has
$E_{-\sigma}$, the first match transfers it to $q_3$. Otherwise
the vanishing of $20$ in the representation and nonvanishing of all
$20$ flags transfer $\mathrm{InV}$ to $q_3$, using the reflected
$31$ vanishing. Both alternatives give ordinary goodness.
\item At $q_4$ the middle $10$ match transfers $E_{-\alpha}$ from
$q_3$, and the subsequent $20$ match transfers $E_{-\sigma}$ from
$q_2$. Their conjunction gives ordinary goodness.
\item At $q_2$, $\mathrm{SameAl}$ with $z$ kills the $10$ and hence
the $32$ representation entries. Removing these from its $Q_0$
support leaves support in the required set for $\mathrm{InA}$;
its $10$ flags are nonzero by the same match. With its
$E_{-\sigma}$ this is the second $\mathrm{Good4}$ clause.
The last point $z$ has the third clause by its seed properties.
\end{itemize}
For each $\mathrm{SamePs}$ match the two pairs lie in $N_P$ with
flags in $Q_0$; for each $\mathrm{SameAl}$ match both representations
have $N_M$ support. Use the corresponding actual polynomial lines
of UX4. For clarity, interpolate the representation and flag matrices
entrywise with parameter $t$. The two endpoints on a $P$ line have
identical diagonal blocks in the ordering $(0,2\mid 1,3)$ (up to
the choice of which block comes first), and their differences occupy
one off-diagonal block. For an $M$ line use the ordering
$(0,1\mid2,3)$; matching $10$ determines the other diagonal block by
the multiplier equation. The differences of the lower flags also
lie in the relevant square-zero block, since they are lower
symplectic unipotents and have matching $20$ or $10$ coordinates.
Products and adjoints therefore give the interpolated endpoint
identities as in (UX.22). In particular conjugation by the interpolated
flag has the prescribed upper shape and diagonal characters. The other
frame conditions follow by evaluating the fixed entries and the
locally constant coefficient functions. This gives four links over
$k[t]$ in the order displayed, with every destination good and the
starting point good. Concatenation gives the literal $\mathrm{Chain4}$
from $u$ to the reduction of the one chosen seed. The construction
of $T_1$ preceded every later $T$, and the seed on $T$ preceded $u$.
This proves $\mathrm{RankFourChainReductionDense}$ under its source
cohomology, fitness and seed hypotheses.
% END RFCR4

\endgroup
